% ==============================================================================
\section{Computational Experiments}
\label{sec:experiments}
% ==============================================================================

This section describes the benchmark design and instance generation procedure.
Detailed numerical results will be presented in subsequent sections once the
full experimental run is complete.

% ------------------------------------------------------------------------------
\subsection{Benchmark Design}
\label{sec:benchmark_design}
% ------------------------------------------------------------------------------

The benchmark is designed to isolate the effect of four axes of variation:
\begin{enumerate}
  \item \textbf{Topology} (six levels): the structure of the blocking-arc graph
        $G = (\calP, \calA)$.
  \item \textbf{Fleet size} (four levels): $R \in \{5, 10, 20, 30\}$ aircraft.
  \item \textbf{Slack} (three levels): how much temporal headroom aircraft have
        between their processing time and their time window.
  \item \textbf{Maximum constraint density}: a hard large-scale configuration
        that combines maximum blocking with a large fleet.
\end{enumerate}

All instances use $P = 5$ hangar positions.  The benchmark comprises
12 configurations, each replicated with 10 independently seeded random
instances, yielding \textbf{120 instances} in total.
Table~\ref{tab:benchmark_matrix} shows the full specification.

\begin{table}[htbp]
  \centering
  \caption{Benchmark specification: 12 configurations $\times$ 10 seeds
           $= 120$ instances.  All instances use $P = 5$ positions.}
  \label{tab:benchmark_matrix}
  \begin{tabular}{llcccl}
    \toprule
    Axis & Topology & $R$ & Slack & Seeds & Purpose \\
    \midrule
    \multirow{6}{*}{Topology}
      & None         & 10 & tight  & 1--10 & Baseline (no blocking) \\
      & Chain        & 10 & tight  & 1--10 & Linear blocking \\
      & Hub          & 10 & tight  & 1--10 & Star topology \\
      & Triangle     & 10 & tight  & 1--10 & Classic hangar layout \\
      & Two rows     & 10 & tight  & 1--10 & Parallel rows \\
      & Full         & 10 & tight  & 1--10 & Maximum blocking \\
    \midrule
    \multirow{3}{*}{Size}
      & Triangle & 5  & tight  & 1--10 & Small (MILP-tractable) \\
      & Triangle & 20 & tight  & 1--10 & Medium (heuristic regime) \\
      & Triangle & 30 & tight  & 1--10 & Large \\
    \midrule
    \multirow{2}{*}{Slack}
      & Triangle & 10 & loose  & 1--10 & Relaxed time windows \\
      & Triangle & 10 & medium & 1--10 & Moderate time windows \\
    \midrule
    Hard large
      & Full & 20 & tight  & 1--10 & Dense + large \\
    \bottomrule
  \end{tabular}
\end{table}

% ------------------------------------------------------------------------------
\subsection{Instance Generation}
\label{sec:instance_gen}
% ------------------------------------------------------------------------------

For each configuration of Table~\ref{tab:benchmark_matrix}, the 10 random
seeds drive a deterministic generator that produces the hangar topology,
the fleet of aircraft, and their associated jobs and time windows.  We
describe each step below.

\subsubsection{Hangar topology families}

Each instance is built on top of one of six canonical blocking-arc graphs
that together span the full spectrum from fully independent to maximally
constrained hangars; this is the source of the topology axis in
Table~\ref{tab:benchmark_matrix}.  In all cases positions are labelled
$\{P1, \dots, P5\}$ in arrival order.  The \emph{None} topology takes
$\calA = \emptyset$, leaving all positions independent and serving as a
baseline with no blocking interactions.  The \emph{Chain} topology
arranges positions in a linear row, $\calA = \{(p_i, p_{i+1}) : 1 \le i < P\}$,
so that each position blocks the next; this represents a single-file
hangar.  The \emph{Hub} topology designates an \emph{aisle} position $p_1$
that blocks all others, $\calA = \{(p_1, p_k) : k \ne 1\}$, modelling a
star layout in which one slot guards the access to the rest.  The
\emph{Triangle} topology turns the last three positions into a transitive
triangle, $\calA = \{(a, b), (a, c), (b, c)\}$ with $\{a, b, c\}$ the rear
trio, reproducing the layout of the facility that originally motivated
this study.  The \emph{Two rows} topology splits positions into a front
and a rear half and pairs them position-by-position, modelling parallel
rows in which only the front aircraft obstructs its corresponding rear
one.  Finally, the \emph{Full} topology takes every ordered pair as a
blocking arc, $\calA = \{(p_i, p_j) : i < j\}$, providing the maximally
constrained extreme.

These six families were chosen so that each one isolates a distinct
qualitative feature of the blocking graph---no blocking, sequential
blocking, star blocking, transitive blocking, paired blocking, and dense
blocking---while keeping the number of positions fixed at $P = 5$.  The
resulting topology axis is therefore a controlled probe of how the
blocking structure, independently of fleet size and time-window tightness,
influences the difficulty of the AHPP and the relative performance of the
solving methods.

\subsubsection{Aircraft and job generation}

Processing times, earliest starts and due dates are sampled on an
integer-day grid; the only fractional quantity in the model is the
tow-in/tow-out time $\varepsilon$ (\S\ref{sec:problem}), which is recorded
in each instance file and fixed at $\varepsilon = 0.5$ days throughout
the reported benchmark.  No \emph{operational} horizon is imposed: only
makespan and delay enter the objective, so any feasible schedule can be
expressed in finite time without an operational upper bound.  The MILP
formulation does use a technical upper bound $H^{UB}$ (defined in
\S\ref{sec:milp}) solely to bound the time variables and calibrate the
big-$M$ constants, but $H^{UB}$ is derived from the instance data rather
than fixed at instance generation, so it is a property of the solver
model and not of the benchmark itself.

Given a configuration with $R$ aircraft and a task-count range
$[n_{\min}, n_{\max}]$, each aircraft $r$ is generated as follows.  Its
number of jobs is drawn as
$n_r \sim \mathrm{Uniform}\{n_{\min}, \ldots, n_{\max}\}$ and the
processing time of each job is sampled independently from the discrete
uniform $D_j \sim \mathrm{Uniform}\{3, 4, 5, 6, 7, 8\}$ days (mean $5.5$,
standard deviation $\approx 1.7$).  The total processing time of the
aircraft is the integer sum
$D_r = \sum_{j \in \calJ_r} D_j$ days.  The earliest start is sampled
uniformly as $E_r \sim \mathrm{Uniform}\{0, \ldots, S\}$ days with arrival
spread $S = \lceil R\,\bar p/P\rceil$ and $\bar p = 5.5$; this scales
naturally with the fleet size and the number of positions without
introducing an operational horizon.  Given the slack ratio $\rho$ of
Table~\ref{tab:slack}, the target finish is set to
$L_r = E_r + D_r + \max\!\bigl(1,\,\lceil \rho\, D_r\rceil\bigr)$ days,
which guarantees a strictly positive integer-day margin over the
processing block.  Each aircraft is labelled with one of
$n_c = \max(2, \min(6, \lfloor R/2 \rfloor))$ client identifiers, which
are kept in the schema to support priority extensions but do not enter
the present objective.  The jobs of each aircraft are linked into a
strict linear precedence chain
$j^r_1 \to j^r_2 \to \cdots \to j^r_{n_r}$, and all random draws are
controlled by the seed so each instance is fully reproducible from its
configuration and seed value.

\begin{table}[htbp]
  \centering
  \caption{Slack level definitions.  $\rho$ is the ratio of post-window
           slack to total aircraft duration $D_r$, applied as
           $L_r = E_r + D_r + \max(1, \lceil \rho D_r \rceil)$.}
  \label{tab:slack}
  \begin{tabular}{lc}
    \toprule
    Level & $\rho$ \\
    \midrule
    Loose  & 0.80 \\
    Medium & 0.35 \\
    Tight  & 0.05 \\
    \bottomrule
  \end{tabular}
\end{table}

% ------------------------------------------------------------------------------
\subsection{Experimental Setup}
\label{sec:exp_setup}
% ------------------------------------------------------------------------------

All MILP and FAS components were modelled with Pyomo~\citep{Bynum2021}
and solved with Gurobi~11~\citep{Gurobi2023}; the topology-aware GRASP
was implemented in plain Python.  Experiments were run on a workstation
with an Intel Core i9 processor (16 cores, 32 threads, 3.5\,GHz base
clock), 32\,GB of RAM, and Windows~11 as the operating system.  Each run
was given a wall-clock time limit of \textbf{60 seconds}.
Table~\ref{tab:methods} lists the four methods evaluated, derived from
the algorithms of \S\ref{sec:approaches}.
Each method is run under three weight profiles:
default $(W^M, W^D, W^S) = (0.1, 1.0, 10)$ (movement-priority),
delay-priority $(1, 10, 0.1)$, and makespan-priority $(10, 0.1, 1)$.
This yields $4 \times 3 = 12$ solver configurations, applied to all 120
instances, for a total of \textbf{1\,440 runs}.

\begin{table}[htbp]
  \centering
  \caption{Solving methods evaluated in the benchmark.}
  \label{tab:methods}
  \begin{tabular}{llll}
    \toprule
    Label & Method & Solver & Notes \\
    \midrule
    \texttt{milp\_baseline}
      & Full MILP
      & Gurobi (via Pyomo)
      & Exact; 60\,s limit \\
    \texttt{topology\_ms6}
      & GRASP, $n_s{=}6$
      & Custom
      & 10\,s per start \\
    \texttt{fas\_on\_topo}
      & FAS on topology assignment
      & Gurobi (direct)
      & 60\,s; fixed $\pi$ \\
    \texttt{safe\_pipeline}
      & Best of \{topology, FAS\}
      & ---
      & Zero extra cost \\
    \bottomrule
  \end{tabular}
\end{table}

% ------------------------------------------------------------------------------
\subsection{Results}
\label{sec:results}
% ------------------------------------------------------------------------------

We report mean values over nine random seeds per configuration.
The analysis is organised around four questions: (i) up to what size is the MILP
provably optimal within the time budget, (ii) how the four methods compare on
the resulting feasible solutions, (iii) how the gap between the heuristic
pipeline and the exact solver evolves across instance types, and
(iv) how the chosen weight profile shapes the trade-off between makespan,
delay, and movements.

\subsubsection{Reach of the exact MILP}

Table~\ref{tab:res_milp_opt} reports the MILP outcome under the 60\,s
budget.  The exact solver only closes the $R{=}5$ configuration to
optimality (in about $4$\,s on average).  On every $R{=}10$ configuration
the search exhausts the budget with an average Gurobi gap above 95\,\%
--- the gap remains high even on the unconstrained \emph{None} topology
($96.2\,\%$), indicating that bound weakness, rather than blocking
interactions, is what stalls branch-and-bound.  On the medium-size
configurations ($R{=}20$, both triangle and full topologies) the solver
overruns the 60\,s wall clock by a factor of two to three (an artefact
of Gurobi's post-timeout cleanup) and still ends with a gap close to
$100\,\%$; moreover, on most seeds the solver either aborts or exceeds the
Pyomo build cap before returning any feasible incumbent: across
nine seeds, Triangle $R{=}20$ yields only 2 usable MILP runs and
Full $R{=}20$ yields just 1.  At $R{=}30$
the model either aborts or exceeds the one-minute Pyomo build cap on
every seed, so no MILP row can be reported for that size; the heuristic
pipeline is therefore the only available source of solutions for the
largest configuration.

% This file is autogenerated by scripts/tests/aggregate_results.py.
% Do not edit by hand — rerun the script to regenerate.
\begin{table}[htbp]
  \centering\small
  \caption{MILP performance under the 60\,s budget per configuration. $n$ = number of usable runs (out of 10 seeds; rows aborted by Gurobi or whose Pyomo build exceeded 60\,s are excluded); \textit{Opt} = number of instances proved optimal; $\overline{\mathrm{gap}}$ = mean Gurobi optimality gap; $\overline{t_{\mathrm{MILP}}}$ = mean wall-clock time (s).}
  \label{tab:res_milp_opt}
  \begin{tabular}{lrrrr}
    \toprule
    Configuration & $n$ & Opt & $\overline{\mathrm{gap}}$ (\%) & $\overline{t_{\mathrm{MILP}}}$ \\
    \midrule
    None ($R{=}10$, tight) & 10 & 1 & 86.36 & 62.1 \\
    Chain ($R{=}10$, tight) & 10 & 0 & 97.73 & 67.3 \\
    Hub ($R{=}10$, tight) & 10 & 0 & 93.89 & 74.5 \\
    Triangle ($R{=}10$, tight) & 10 & 0 & 95.03 & 67.7 \\
    Two rows ($R{=}10$, tight) & 10 & 0 & 94.69 & 75.1 \\
    Full ($R{=}10$, tight) & 10 & 0 & 98.44 & 67.0 \\
    \midrule
    Triangle ($R{=}5$, tight) & 10 & 10 & 0.00 & 5.0 \\
    Triangle ($R{=}20$, tight) & 3 & 0 & 99.67 & 130.9 \\
    Triangle ($R{=}30$, tight) & 1 & 0 & 99.71 & 61.5 \\
    \midrule
    Triangle ($R{=}10$, medium) & 10 & 0 & 92.31 & 68.7 \\
    Triangle ($R{=}10$, loose) & 10 & 1 & 55.88 & 61.9 \\
    \midrule
    Full ($R{=}20$, tight) & 2 & 0 & 99.81 & 109.0 \\
    \bottomrule
  \end{tabular}
\end{table}


\subsubsection{Per-method means by configuration}

Table~\ref{tab:res_default} aggregates the four methods under the
default profile $(W^M, W^D, W^S) = (0.1, 1, 10)$.  Three observations
stand out.  First, on every $R{=}10$ configuration with non-trivial
blocking the heuristic pipeline finishes well below the best feasible
MILP solution found within the budget: \textbf{Safe} reduces $\bar f$ by
$20$--$40\,\%$ over \textbf{MILP} on Chain, Hub, Triangle, Two rows and
Full, while on the unconstrained \emph{None} topology the two methods
are practically tied ($104.6$ versus $106.3$).  Within a 60\,s budget
the heuristic therefore dominates the exact solver on the constrained
topologies and matches it on the trivial one.  Second, the gap widens
sharply with the fleet size: on \emph{Triangle $R{=}20$} \textbf{Safe}
returns $\bar f = 678$ against \textbf{MILP}'s $3250$, and on
\emph{Full $R{=}20$} the gap is $956$ versus $4676$, both more than
four times apart.  Third, the topology heuristic alone (\textbf{Topo})
already matches the safe-pipeline value on most $R{=}10$ rows; on
\emph{Chain}, \emph{Two rows}, \emph{Hub} and \emph{Full} a single FAS
pass shaves a few percent off Topo, and on \emph{Triangle $R{=}30$}
\textbf{Safe} ends below both Topo and FAS in the mean because it
picks the per-seed minimum --- this is by construction and illustrates
the value of the safe-pipeline composition.

% This file is autogenerated by scripts/tests/aggregate_results.py.
% Do not edit by hand — rerun the script to regenerate.
\begin{table}[htbp]
  \centering
  \caption{Mean results over 6 seeds per configuration under the default weight profile $(W^M, W^D, W^S) = (0.1, 1, 10)$. Columns per method: mean objective $\bar f$, mean makespan $\bar m$, mean total delay $\bar v^D$, mean movements $\bar n$, and mean wall-clock time $\bar t$ (s). MILP is unavailable for $R=30$.}
  \label{tab:res_default}
  \setlength{\tabcolsep}{3pt}%
  \resizebox{\textwidth}{!}{%
  \begin{tabular}{lrrrrrrrrrrrrrrrrrrrr}
    \toprule
    Configuration & \multicolumn{5}{c}{\textbf{MILP}} & \multicolumn{5}{c}{\textbf{Topo}} & \multicolumn{5}{c}{\textbf{FAS}} & \multicolumn{5}{c}{\textbf{Safe}} \\
     & $\bar f$ & $\bar m$ & $\bar v^D$ & $\bar n$ & $\bar t$ & $\bar f$ & $\bar m$ & $\bar v^D$ & $\bar n$ & $\bar t$ & $\bar f$ & $\bar m$ & $\bar v^D$ & $\bar n$ & $\bar t$ & $\bar f$ & $\bar m$ & $\bar v^D$ & $\bar n$ & $\bar t$ \\
    \midrule
    None ($R{=}10$, tight) & 109.28 & 66.17 & 102.67 & 0.0 & 69.1 & 107.22 & 65.50 & 100.67 & 0.0 & 60.0 & 107.22 & 65.50 & 100.67 & 0.0 & 0.0 & 107.22 & 65.50 & 100.67 & 0.0 & 0.0 \\
    Chain ($R{=}10$, tight) & 220.17 & 90.00 & 197.83 & 1.3 & 69.4 & 177.32 & 85.67 & 158.75 & 1.0 & 60.0 & 174.12 & 86.17 & 155.50 & 1.0 & 0.2 & 174.12 & 86.17 & 155.50 & 1.0 & 0.0 \\
    Hub ($R{=}10$, tight) & 197.60 & 90.17 & 188.58 & 0.0 & 80.9 & 142.16 & 82.42 & 133.92 & 0.0 & 60.0 & 142.03 & 82.75 & 133.75 & 0.0 & 0.1 & 142.03 & 82.75 & 133.75 & 0.0 & 0.0 \\
    Triangle ($R{=}10$, tight) & 145.82 & 79.08 & 137.92 & 0.0 & 68.3 & 116.78 & 69.50 & 109.83 & 0.0 & 60.0 & 116.78 & 69.50 & 109.83 & 0.0 & 0.0 & 116.78 & 69.50 & 109.83 & 0.0 & 0.0 \\
    Two rows ($R{=}10$, tight) & 146.71 & 77.92 & 138.92 & 0.0 & 81.9 & 116.15 & 68.17 & 109.33 & 0.0 & 60.1 & 113.08 & 67.50 & 106.33 & 0.0 & 0.4 & 113.08 & 67.50 & 106.33 & 0.0 & 0.0 \\
    Full ($R{=}10$, tight) & 303.73 & 119.75 & 275.08 & 1.7 & 68.6 & 195.76 & 86.75 & 180.42 & 0.7 & 60.0 & 185.63 & 86.33 & 173.67 & 0.3 & 0.4 & 185.63 & 86.33 & 173.67 & 0.3 & 0.0 \\
    \midrule
    Triangle ($R{=}5$, tight) & 3.55 & 35.50 & 0.00 & 0.0 & 4.9 & 3.72 & 35.50 & 0.17 & 0.0 & 60.1 & 3.72 & 35.50 & 0.17 & 0.0 & 0.4 & 3.72 & 35.50 & 0.17 & 0.0 & 0.0 \\
    Triangle ($R{=}20$, tight) & 2425.10 & 321.00 & 2393.00 & 0.0 & 194.4 & 687.12 & 132.00 & 673.92 & 0.0 & 63.1 & 677.12 & 132.00 & 663.92 & 0.0 & 23.4 & 677.12 & 132.00 & 663.92 & 0.0 & 0.0 \\
    Triangle ($R{=}30$, tight) & --- & --- & --- & --- & --- & 2140.12 & 231.25 & 2107.00 & 1.0 & 588.0 & 2146.20 & 234.50 & 2092.75 & 3.0 & 60.5 & 2119.82 & 234.08 & 2079.75 & 1.7 & 0.0 \\
    \midrule
    Triangle ($R{=}10$, medium) & 140.62 & 94.58 & 131.17 & 0.0 & 70.3 & 74.38 & 68.83 & 67.50 & 0.0 & 60.0 & 74.38 & 68.83 & 67.50 & 0.0 & 0.1 & 74.38 & 68.83 & 67.50 & 0.0 & 0.0 \\
    Triangle ($R{=}10$, loose) & 17.50 & 70.83 & 10.42 & 0.0 & 68.5 & 15.14 & 71.42 & 8.00 & 0.0 & 60.0 & 15.14 & 71.42 & 8.00 & 0.0 & 0.1 & 15.14 & 71.42 & 8.00 & 0.0 & 0.0 \\
    \midrule
    Full ($R{=}20$, tight) & 4675.60 & 531.00 & 4622.50 & 0.0 & 156.2 & 992.38 & 167.08 & 962.33 & 1.3 & 62.8 & 972.98 & 168.17 & 922.83 & 3.3 & 65.6 & 963.27 & 170.25 & 926.25 & 2.0 & 0.0 \\
    \bottomrule
  \end{tabular}}
\end{table}


\subsubsection{Relative gap to the best solution}

To compare methods on a common scale we compute, for each instance, the
\emph{relative gap} of each method to the best objective achieved by the
four methods on that instance.  The gaps in
Table~\ref{tab:res_method_gap} confirm that, within the 60\,s budget,
the exact solver is the worst performer on every configuration with
non-trivial blocking: from $22\,\%$ on \emph{Triangle $R{=}10$} up to
$47\,\%$ on \emph{Hub $R{=}10$} and $68\,\%$ on \emph{Full $R{=}10$},
and dramatic gaps on the larger instances ($417\,\%$ on
\emph{Triangle $R{=}20$}, $431\,\%$ on \emph{Full $R{=}20$}), driven by
the very weak feasible incumbents Gurobi returns at the time limit.
\textbf{Topo} alone stays below $7\,\%$ on every $R{=}10$ tight
configuration and below $4\,\%$ at $R{=}20$; \textbf{FAS} closes the
residual gap to $0\,\%$ in every case except \emph{Triangle $R{=}30$},
where Topo and FAS each win on some seeds and only the Safe pipeline
attains $0\,\%$ throughout.  Two apparent outliers stand out: the heuristic methods sit at
$32.2\,\%$ on \emph{Triangle $R{=}10$, loose} (MILP at $61.5\,\%$),
and at $3.5\,\%$ on \emph{Triangle $R{=}5$} (Topo/FAS/Safe).  Both come from
seeds whose best objective is close to zero, so even tiny absolute
deviations translate to large relative gaps; the absolute differences
in Table~\ref{tab:res_default} remain well within a few units.

% This file is autogenerated by scripts/tests/aggregate_results.py.
% Do not edit by hand — rerun the script to regenerate.
\begin{table}[htbp]
  \centering\small
  \caption{Mean relative gap (\%) of each method with respect to the best objective found across the four methods on the same seed, under the default weight profile. The Safe pipeline is the per-seed minimum of Topo and FAS, so its gap is always $\leq$ both. Lower is better.}
  \label{tab:res_method_gap}
  \begin{tabular}{lrrrr}
    \toprule
    Configuration & MILP & Topo & FAS & Safe \\
    \midrule
    None ($R{=}10$, tight) & 1.60 & 0.06 & 0.06 & 0.06 \\
    Chain ($R{=}10$, tight) & 23.03 & 1.31 & 0.00 & 0.00 \\
    Hub ($R{=}10$, tight) & 47.44 & 0.15 & 0.09 & 0.09 \\
    Triangle ($R{=}10$, tight) & 22.10 & 0.00 & 0.00 & 0.00 \\
    Two rows ($R{=}10$, tight) & 34.66 & 1.71 & 0.05 & 0.05 \\
    Full ($R{=}10$, tight) & 67.70 & 4.54 & 0.00 & 0.00 \\
    \midrule
    Triangle ($R{=}5$, tight) & 0.00 & 3.47 & 3.47 & 3.47 \\
    Triangle ($R{=}20$, tight) & 416.81 & 1.36 & 0.00 & 0.00 \\
    Triangle ($R{=}30$, tight) & --- & 1.24 & 1.88 & 0.00 \\
    \midrule
    Triangle ($R{=}10$, medium) & 65.29 & 0.00 & 0.00 & 0.00 \\
    Triangle ($R{=}10$, loose) & 61.48 & 32.23 & 32.23 & 32.23 \\
    \midrule
    Full ($R{=}20$, tight) & 430.87 & 3.56 & 0.65 & 0.00 \\
    \bottomrule
  \end{tabular}
\end{table}


\subsubsection{Weight-profile sensitivity}

Table~\ref{tab:res_weights} isolates the effect of the three weight
profiles on the structure of the safe-pipeline solutions.  Under the
default, movement-priority profile, the safe pipeline returns
essentially zero blocking movements on every configuration covered by
the benchmark --- the only non-zero entries are residuals of
$\bar n{=}0.9$ on \emph{Chain $R{=}10$}, $\bar n{=}0.2$ on
\emph{Full $R{=}10$}, $\bar n{=}2.2$ on \emph{Triangle $R{=}30$} and
$\bar n{=}2.0$ on the dense \emph{Full $R{=}20$} --- at the cost of a
moderate makespan penalty: e.g.\ $\bar m{=}89.5$ on
\emph{Full $R{=}10$} versus $62.6$ under wC, and $\bar m{=}166.7$ on
\emph{Full $R{=}20$} versus $118.6$.  Switching to the delay-priority
profile (wB) collapses the mean delay onto a narrow band ---
$\bar v^D \approx 98.1$ on every $R{=}10$ tight row regardless of
topology --- but injects $6$--$44$ blocking movements per instance.
The makespan-priority profile (wC) flattens $\bar m$ to about
$62.5$ on every $R{=}10$ tight row at the price of
$4$--$30$ movements and intermediate delay.  The three profiles
therefore act as effective dials between the objectives, and the
default profile is the appropriate choice when the operational
priority is to keep aircraft in the assigned slot for the whole
servicing window.

% This file is autogenerated by scripts/tests/aggregate_results.py.
% Do not edit by hand — rerun the script to regenerate.
\begin{table}[htbp]
  \centering\small
  \caption{Effect of the weight profile on the Safe pipeline solution structure. Mean $\bar m$, mean total delay $\bar v^D$, and mean movements $\bar n$ (averaged over 9 seeds) under the three profiles: default $(0.1,1,10)$, wB $(1,10,0.1)$ (delay-priority), wC $(10,0.1,1)$ (makespan-priority).}
  \label{tab:res_weights}
  \resizebox{\textwidth}{!}{%
  \begin{tabular}{lrrrrrrrrr}
    \toprule
    Configuration & \multicolumn{3}{c}{\textbf{Default}} & \multicolumn{3}{c}{\textbf{wB (delay)}} & \multicolumn{3}{c}{\textbf{wC (makespan)}} \\
     & $\bar m$ & $\bar v^D$ & $\bar n$ & $\bar m$ & $\bar v^D$ & $\bar n$ & $\bar m$ & $\bar v^D$ & $\bar n$ \\
    \midrule
    None ($R{=}10$, tight) & 65.39 & 98.06 & 0.0 & 64.83 & 98.06 & 0.0 & 62.50 & 102.28 & 0.0 \\
    Chain ($R{=}10$, tight) & 88.17 & 153.78 & 0.9 & 64.83 & 98.06 & 17.8 & 62.50 & 110.06 & 11.3 \\
    Hub ($R{=}10$, tight) & 78.72 & 128.06 & 0.0 & 65.39 & 98.06 & 16.4 & 62.50 & 115.17 & 11.3 \\
    Triangle ($R{=}10$, tight) & 68.61 & 106.94 & 0.0 & 64.83 & 98.06 & 10.9 & 62.50 & 106.39 & 8.0 \\
    Two rows ($R{=}10$, tight) & 67.72 & 105.39 & 0.0 & 65.39 & 98.06 & 5.8 & 62.50 & 106.33 & 3.8 \\
    Full ($R{=}10$, tight) & 89.50 & 181.22 & 0.2 & 65.39 & 98.06 & 44.0 & 62.61 & 119.17 & 29.6 \\
    \midrule
    Triangle ($R{=}5$, tight) & 35.89 & 0.11 & 0.0 & 35.89 & 0.00 & 0.2 & 35.89 & 3.11 & 0.0 \\
    Triangle ($R{=}20$, tight) & 132.50 & 662.72 & 0.2 & 123.56 & 616.67 & 33.6 & 117.50 & 645.67 & 23.6 \\
    Triangle ($R{=}30$, tight) & 234.72 & 2060.78 & 2.2 & 219.50 & 1933.56 & 60.5 & 207.25 & 2003.56 & 42.2 \\
    \midrule
    Triangle ($R{=}10$, medium) & 68.17 & 63.06 & 0.0 & 64.94 & 51.06 & 10.9 & 62.50 & 62.28 & 7.6 \\
    Triangle ($R{=}10$, loose) & 70.56 & 5.78 & 0.0 & 67.06 & 0.33 & 9.6 & 62.50 & 11.56 & 8.9 \\
    \midrule
    Full ($R{=}20$, tight) & 166.72 & 919.17 & 2.0 & 123.78 & 618.72 & 123.8 & 118.61 & 681.56 & 84.4 \\
    \bottomrule
  \end{tabular}}
\end{table}


\FloatBarrier
